% !TEX root = ../main.tex
\chapter{全书路线图、符号、坐标与基本约定}
\label{ch:roadmap}

\section*{学习目标}
\begin{itemize}
  \item 建立本书统一的坐标系、时间因子与频率约定。
  \item 统一“模、Bloch 模、$\Gamma$ 点、带边、被动共振、阈值增益、阈值电流”等术语含义。
  \item 明确本书中不同模型层级的关系，避免跨层误读结论。
\end{itemize}

\section*{先修知识}
复数表示、向量分析、傅里叶级数与基本晶体学记号。

\section*{本章路线图}
\ChapterModelLevel{L1→L4（层级接口总览）}
本章不推导新的器件结果，而是固定全书使用的共同语言：先固定 benchmark 教学器件与高频符号工具页，再统一坐标和频域约定，随后说明模型层级与术语边界，最后给出阅读全书时必须守住的六条“红线”。晶格与倒格矢记号在 \cref{ch:bloch} 首用处建立，归一化约定在 \cref{ch:cwt-min} 首用处说明。

\section{为什么在第二章先停下来统一语言}
初学者读 PCSEL 文献时，最先遇到的困难往往不是公式本身，而是同一个词在不同上下文里意思变了。比如“阈值”这个词，在某些地方指的是阈值增益，在另一些地方指的是阈值电流；“模式”有时指无限周期的 Bloch 模，有时又指有限器件里的实际振荡模；“TE 模”在二维教材里是严格分类，在 slab 结构里却经常只是近似称呼。若不先限定这些边界，后续公式和结论容易被跨层误读。

因此，本章的作用不是提供一份可有可无的记号表，而是给后续所有推导建立语义基础。它相当于为全书固定坐标、符号和模型层级。

\section{本书默认的 benchmark 教学器件}
为了避免读者在前半本不断在脑中切换不同器件，本书从这里开始固定一个默认教学参照：\textbf{方格晶格 air-hole slab PCSEL}。除非某章明确声明换了结构，否则本书的几何、对称性、PWEM 候选模、CWT 四波图像、单胞筛选、有限阵列检查与 worked example 都默认围绕这一 benchmark 展开。

它的最小身份是：
\begin{itemize}
  \item 平面晶格：方格晶格；
  \item 单胞类型：高折射率背景中的圆形空气孔；
  \item 纵向结构：有源 slab 波导，上下包层不必严格对称；
  \item 材料平台：以 GaAs-like 近红外半导体平台作教学参照；
  \item 目标波长：近红外 \SIrange{0.8}{1.0}{\micro\meter} 量级；
  \item 目标问题：从 $\Gamma$ 点候选模出发，逐层推进到有限尺寸、电注入与热回写。
\end{itemize}

\begin{figure}[htbp]
  \centering
  % Benchmark device schematic used across the book
\begin{tikzpicture}[x=1cm,y=1cm,>=Latex]
  % (a) Top view
  \begin{scope}[shift={(0,0)}]
    \node[anchor=south,font=\bfseries] at (2.2,3.0) {(a) 俯视：方格孔阵};
    \draw[thick] (0,0) rectangle (4.4,2.7);
    \foreach \x in {0.5,1.3,2.1,2.9,3.7} {
      \foreach \y in {0.45,1.15,1.85,2.55} {
        \draw[fill=white] (\x,\y) circle (0.16);
      }
    }
    \draw[<->,thick] (0.5,-0.25)--(1.3,-0.25);
    \node[below] at (0.9,-0.25) {$a$};
    \draw[<->,thick] (2.1,1.15)--(2.26,1.15);
    \node[above right] at (2.26,1.15) {$r$};
    \node[below right] at (0,2.7) {周期边界方向 $x,y$};
  \end{scope}

  % (b) Side view
  \begin{scope}[shift={(6.2,0)}]
    \node[anchor=south,font=\bfseries] at (2.7,3.0) {(b) 侧视：外延层栈};
    \draw[fill=purple!15] (0,1.7) rectangle (5.4,2.35);
    \draw[fill=red!18]    (0,1.25) rectangle (5.4,1.7);
    \draw[fill=green!14]  (0,0.9) rectangle (5.4,1.25);
    \draw[fill=blue!12]   (0,0.0) rectangle (5.4,0.9);
    \foreach \x in {0.45,1.35,2.25,3.15,4.05,4.95} {
      \draw[fill=white] (\x,2.02) circle (0.16);
    }
    \draw[thick] (0,0) rectangle (5.4,2.35);

    \draw[<->] (5.65,1.7)--(5.65,2.35); \node[right] at (5.65,2.02) {$t_{\mathrm{pc}}$};
    \draw[<->] (5.65,1.25)--(5.65,1.7); \node[right] at (5.65,1.48) {$t_{\mathrm{core}}$};
    \draw[<->] (5.65,0.9)--(5.65,1.25); \node[right] at (5.65,1.07) {$t_{\mathrm{clad}}$};
    \draw[<->] (5.65,0.0)--(5.65,0.9); \node[right] at (5.65,0.45) {$t_{\mathrm{sub}}$};

    \node[left] at (0,2.02) {\scriptsize PC 层};
    \node[left] at (0,1.47) {\scriptsize 波导芯层/MQW};
    \node[left] at (0,1.08) {\scriptsize 包层};
    \node[left] at (0,0.45) {\scriptsize 衬底};
  \end{scope}
\end{tikzpicture}


  \caption{本书统一教学器件示意。(a) 俯视图给出方格孔阵和几何参数 $a,r$；(b) 侧视图给出层栈与常用纵向参数（$t_{\mathrm{pc}},t_{\mathrm{core}},t_{\mathrm{sub}}$ 等）。后续章节若未特别声明，默认都以此结构作为教学对象。}
  \label{fig:ch02_benchmark_device}
\end{figure}

这并不是说全书结论只对这一类器件成立，而是说：\textbf{若某章没有特别声明，我们优先把抽象公式和分类工具投影回这一个统一对象。} 这样做的好处是，读者从 \cref{ch:bloch,ch:gamma,ch:pwe} 读到 \cref{ch:workflow,ch:worked-example} 时，脑中追踪的是同一个器件在不同层级的投影，而不是每章换一个新结构。

\section{前置工具页：全书最高频符号先一次看全}
为了避免读者在后文反复翻页找定义，本节把全书最高频、且最容易跨层混淆的量先压成一张工具表。更完整版本见书末术语表，但这里先给出读书时最常需要随手回看的核心对象。

\begin{table}[htbp]
  \centering
  \scriptsize
  \caption{全书最高频核心符号的前置工具页。表中“典型量级”用于\textbf{[教学示意]}与\textbf{[本书算例]}对齐，跨平台定量比较需回到各章文献值与边界条件。}
  \label{tab:core_symbols_toolpage}
  \setlength{\tabcolsep}{3pt}
  \begin{tabularx}{\textwidth}{@{}p{1.95cm}p{1.95cm}p{2.2cm}Y@{}}
    \toprule
    符号 & 所属层级 & 典型单位/量级 & 最小含义 \\
    \midrule
    $\tilde{\omega}=\omega_r-\ii\gamma$ & L1 & \si{rad.s^{-1}}（$\omega_r\sim10^{15}$） & 开放系统复频率；实部给出中心频率，虚部给出衰减率 \\
    $Q,\tau_p$ & L1/L2 & $Q$ 无量纲；$\tau_p$ 常为 ps--ns & 品质因子与光子寿命；描述被动或线性工作点下的损耗时间尺度 \\
    $g_{\mathrm{mat}},g_{\mathrm{mod}},g_{\mathrm{th}}$ & L2 & \si{cm^{-1}}（常见 $10^1$--$10^3$） & 材料增益、模增益、阈值增益；回答“增益补偿损耗需要到哪一步” \\
    $I_{\mathrm{th}}$ & L3/L4 & \si{mA}--\si{A}（随面积变化） & 阈值电流；真实器件达到阈值工作点所需的注入电流 \\
    $\Gamma_{\mathrm{opt}},\Gamma_{\mathrm{loss}}$ & L1/L2 & 无量纲（常见 $10^{-2}$--$10^{-1}$） & 目标模与有源区、损耗区的重叠因子 \\
    $\alpha_{\mathrm{rad}},\alpha_{\mathrm{int}}$ & L1/L2 & \si{cm^{-1}} 或 \si{m^{-1}} & 辐射损耗与内部损耗的等效损耗系数 \\
    $N,S,T,J$ & L3/L4 & $N$:\si{cm^{-3}}；$T$:\si{K}；$J$:\si{A.cm^{-2}} & 载流子、光子数、温度、电流密度；器件动态与多物理场的最小状态量 \\
    $\E(\rvec)$、角谱 $\tilde{E}$ & L1 & 归一化或 \si{V.m^{-1}} & 近场与角谱/远场表示，用于连接模场、偏振与远场 \\
    \bottomrule
  \end{tabularx}
\end{table}

\begin{RigorousBox}
从阅读策略上，本书的默认做法是：先把 benchmark 结构当作统一对象，再把每一章理解为“对同一对象做一次新的投影”。因此，若某章给出的是 L1 被动候选模目录，就不要在脑中把它误升级成 L4 器件工作结论；它只是 benchmark 在当前层级的一个中间表示。
\end{RigorousBox}

\section{空间坐标与结构坐标：先分清什么是纵向，什么是面内}
本书统一采用直角坐标 $\rvec=(x,y,z)$。其中：
\begin{itemize}
  \item $z$ 轴沿外延层法向，也就是器件纵向；
  \item $(x,y)$ 平面是光子晶体的周期平面，也叫面内平面；
  \item 当需要把平面内向量单独拿出来时，记作 $\rho=(x,y)$。
\end{itemize}

这个约定看似简单，但会反复用到。因为对PCSEL而言，纵向与横向承担的是不同物理角色：纵向层栈决定导模约束、外延不对称和吸收层分布；横向周期结构决定Bloch模分类、布拉格耦合与远场工程。后文如果不特别说明，看到"纵向"就默认与外延层栈相关，看到"面内"就默认与二维周期和横向模式选择相关。


\section{时间因子与复频率：为什么符号约定会影响物理解释}
本书统一采用时间因子
\begin{equation}
\ee^{-\ii\omega t}.
\label{eq:time_factor}
\end{equation}
这意味着如果一个本征频率写成
\[
\tilde{\omega}=\omega_0-\ii\gamma,
\]
那么对应场随时间的包络衰减为
\[
\ee^{-\ii\tilde{\omega}t}=\ee^{-\ii\omega_0 t}\ee^{-\gamma t}.
\]
因此在本书的约定下，$\gamma>0$ 代表衰减。这个约定必须记住，因为不同论文可能采用 $\ee^{+\ii\omega t}$，那时复频率虚部的符号会相反。

\subsection{角频率、普通频率与波长}
除非特别说明，本书默认：
\[
\omega=2\pi f,\qquad
k_0=\frac{\omega}{c}=\frac{2\pi}{\lambda}.
\]
其中 $\omega$ 是角频率，$f$ 是普通频率，$\lambda$ 是真空波长。若材料色散明显，折射率、介电常数与增益都应理解为频率相关量，但为了避免公式被符号淹没，后文经常不把这种依赖逐项写出。

\begin{RigorousBox}
一旦进入开放系统，本征频率通常是复数。复频率虚部不是“数值噪声的副产物”，而是模式辐射或损耗的时间尺度编码。若把开放系统错当成只有实频率的封闭系统，就会直接丢失辐射损耗信息。
\end{RigorousBox}

\section{本书的模型层级：为什么同一个问题会有不同答案}
贯穿全书的一个核心思想是：PCSEL 研究不是单一模型，而是一串层级递进的模型。我们采用四层划分。

\begin{figure}[htbp]
  \centering
  \begin{tikzpicture}[node distance=1.6cm,>=Latex]
    \node[draw,rounded corners,fill=blue!8,minimum width=10cm,minimum height=0.95cm] (l1) {L1 被动电磁层：模式、复频率、辐射损耗、远场趋势};
    \node[draw,rounded corners,fill=green!8,minimum width=10cm,minimum height=0.95cm,below=of l1] (l2) {L2 有源光学层：材料增益、模增益、阈值增益、模式竞争};
    \node[draw,rounded corners,fill=yellow!14,minimum width=10cm,minimum height=0.95cm,below=of l2] (l3) {L3 电学层：Poisson、漂移-扩散、注入效率、电流扩展};
    \node[draw,rounded corners,fill=red!10,minimum width=10cm,minimum height=0.95cm,below=of l3] (l4) {L4 热-电-光耦合层：温升、热漂移、工作窗口、稳定性};
    \draw[->,thick] (l1)--(l2);
    \draw[->,thick] (l2)--(l3);
    \draw[->,thick] (l3)--(l4);
  \end{tikzpicture}
  \caption{本书采用的四层建模结构。越往下越接近真实器件工作状态。}
  \label{fig:model_layers}
\end{figure}

\cref{tab:ch02_model_hierarchy_interface} 进一步给出四层模型的接口规范：每一层的输出成为下一层的输入。实际建模时应当按这张表核对数据传递。

\begin{table}[htbp]
  \centering
  \caption{PCSEL 四层建模体系的接口规范。以“计算阈值电流”为例说明数据如何在各层之间传递。每一层的输出成为下一层的输入，形成完整的证据链。}
  \label{tab:ch02_model_hierarchy_interface}
  % Table: L1-L4 model hierarchy interface specification
% Chapter: ch02 - Mathematical Preparation and Modeling Roadmap
% Description: Detailed mapping of inputs/outputs for each modeling layer with concrete examples

\small
\setlength{\tabcolsep}{3.5pt}
\begin{tabularx}{\textwidth}{@{}p{1.8cm}p{2.8cm}p{2.8cm}p{2.8cm}X@{}}
    \toprule
    \textbf{层级} & \textbf{核心问题} & \textbf{输入参数} & \textbf{输出结果} & \textbf{典型求解器} \\
    \midrule
    \textbf{L1: 被动电磁} & 有哪些候选模？哪些能辐射？ & 
    \begin{itemize}[nosep,leftmargin=*]\item 晶格常数$a$\item 填充因子FF\item 刻蚀深度$h$\item 折射率$n(x,y,z)$\end{itemize} & 
    \begin{itemize}[nosep,leftmargin=*]\item 复频率$\tilde{\omega}_n$\item 品质因子$Q_n$\item 辐射损耗$\alpha_{\mathrm{rad}}$\item 远场图案$E(\theta,\phi)$\end{itemize} & 
    PWEM (能带), RCWA (辐射), FDTD (瞬态), FEM (本征模) \\
    \midrule
    \textbf{L2: 有源光学} & 哪个模阈值最低？ & 
    \begin{itemize}[nosep,leftmargin=*]\item L1 输出的$Q_n$\item 内部损耗$\alpha_{\mathrm{int}}$\item 增益谱$g(\lambda, N)$\item 光学限制因子$\Gamma$\end{itemize} & 
    \begin{itemize}[nosep,leftmargin=*]\item 阈值增益$g_{\mathrm{th}}$\item 模增益排序\item 模式竞争结果\item 偏振选择比\end{itemize} & 
    CWT (耦合波), 重叠积分模型，速率方程稳态解 \\
    \midrule
    \textbf{L3: 电注入器件} & 需要多大电流？ & 
    \begin{itemize}[nosep,leftmargin=*]\item L2输出的$g_{\mathrm{th}}$\item 外延掺杂分布\item 电极几何形状\item 迁移率模型$\mu(N,T)$\end{itemize} & 
    \begin{itemize}[nosep,leftmargin=*]\item 阈值电流$I_{\mathrm{th}}$\item 载流子分布$N(x,y,z)$\item 电流密度$J(x,y)$\item 串联电阻$R_s$\end{itemize} & 
    Poisson + 漂移扩散(CHARGE, Semiconductor Module) \\
    \midrule
    \textbf{L4: 热 - 电 - 光耦合} & 能否稳定工作？ & 
    \begin{itemize}[nosep,leftmargin=*]\item L3输出的$P_{\mathrm{heat}}$\item 热导率$\kappa(T)$\item 热边界条件\item dn/dT, dEg/dT\end{itemize} & 
    \begin{itemize}[nosep,leftmargin=*]\item 温度场$T(x,y,z)$\item 热阻$R_{\mathrm{th}}$\item 温升$\Delta T$\item 模式排序变化\end{itemize} & 
    热传导方程(HEAT Module) + 光学回写迭代 \\
    \bottomrule
\end{tabularx}

\vspace{0.5em}
\parbox{\linewidth}{
  \footnotesize
  \textbf{注}: 实际工作流程中，L1→L2→L3→L4 是单向传递，但 L4 产生的温升会通过折射率温度系数($dn/dT \approx 2\times10^{-4}\,\mathrm{K}^{-1}$ for GaAs) 和带隙温度系数($dE_g/dT \approx -0.4\,\mathrm{meV/K}$) 回写到 L1 和 L2，形成闭环耦合。当 $\Delta T > 30\,\mathrm{K}$ 时，这种反馈不可忽略。
}

\end{table}


这四层的区别，不是“谁更高级”，而是“谁回答的问题不同”。

\begin{itemize}
  \item L1 被动电磁层：回答候选模是什么、频率在哪里、辐射损耗如何；
  \item L2 有源光学层：回答要补多少增益才能达到阈值、哪些模式更可能竞争；
  \item L3 电学层：回答注入是否均匀、载流子是否能有效进入有源区；
  \item L4 热-电-光层：回答真实器件在工作时会不会因为温升和非均匀性而改变模式排序。
\end{itemize}

\begin{IntuitionBox}
判断一个具体分析属于哪一层的最直接方法，是看它回答的问题：
\begin{center}
\begin{tabular}{@{}p{6.4cm}p{2.8cm}@{}}
\toprule
若要回答的问题是…… & 最低需要到哪层 \\
\midrule
候选模频率、辐射损耗、远场对称性趋势 & L1 \\
哪个候选模阈值增益最低 & L2 \\
器件阈值电流是多少 & L3 \\
高温下哪个模先激射、线宽如何变化 & L4 \\
\bottomrule
\end{tabular}
\end{center}
若结论跨越了当前层级（例如用被动 $Q$ 直接报告阈值电流），就是“层级偷换”——这是全书反复提醒的最常见错误。
\end{IntuitionBox}

\section{几个最容易混淆的术语：先立边界，细节由首用章展开}
下面四对术语会反复出现。它们的完整辨析放在各自的首用章节（“模 vs Bloch 模”见 \cref{ch:bloch}，“$\Gamma$ 点 vs 带边”见 \cref{ch:gamma}，“被动共振 vs 有源阈值”与“阈值增益 vs 阈值电流”见 \cref{ch:threshold}），高频混淆对照表统一收在 \cref{ch:glossary}。这里只先立四条最低边界：
\begin{itemize}
  \item “模”泛指满足给定方程和边界条件的场分布；“Bloch 模”专指无限周期结构中满足 Bloch 条件的本征场，天然属于 L1 层；
  \item $\Gamma$ 点是布里渊区中心（$\kvec=\bm{0}$）；带边是群速度趋近于零的位置——两者可能重合，但不是同义词；
  \item 被动共振提供候选模，有源阈值决定谁先赢，这个顺序不能反过来；
  \item 阈值增益 $g_\mathrm{th}$ 是 L2 层的光学条件，阈值电流 $I_\mathrm{th}$ 是 L4 层的器件条件，中间隔着电学与热学两层桥。
\end{itemize}

\begin{PitfallBox}
“某模式在 RCWA/FEM 中 $Q$ 最高，因此器件阈值电流最低”这句话同时混淆了被动与有源、阈值增益与阈值电流、无限周期与有限器件三个层级。
\end{PitfallBox}

\section{如何阅读后续章节：每一章都在回答链条上的哪一环}
为了让整本书读起来像一条连续链，而不是二十多篇松散讲义，这里先给出主线。

\begin{equation}
\begin{aligned}
\{\varepsilon(\rvec),\mu(\rvec)\}
&\xrightarrow{\text{Maxwell + Bloch}}
\{\tilde{\omega}_m,\bm{F}_m\}
\\
&\xrightarrow{\text{gain/loss overlap}}
g_{\mathrm{th},m}
\xrightarrow{\text{DD + Poisson + Heat}}
I_{\mathrm{th},m}.
\end{aligned}
\label{eq:book_chain}
\end{equation}

这条链里每个箭头对应后续一批章节：
\begin{itemize}
  \item \cref{ch:maxwell,ch:bloch,ch:gamma} 建立 L1 层；
  \item \cref{ch:feedback,ch:threshold,ch:cwt-min} 解释模耦合与阈值增益；
  \item \cref{ch:epitaxial-optics,ch:qw-gain,ch:electrical} 把外延层和电学结构引入；
  \item \cref{ch:eto-coupling,ch:workflow} 则把这些模块连成真实器件工作流。
\end{itemize}

理解 \cref{eq:book_chain} 的意义在于：后面任何局部公式都不再是孤立的。它们总是在为某一个箭头服务。

\section{全书模型依赖图：每章向下一章交付什么}
仅知道“后面有哪些章节”还不够，还需要知道每一段向下一段交付什么中间量。对教材读者来说，这张依赖表比一串目录更实用，因为它直接说明本章在后续推导中的作用。\cref{tab:chapter-handoff} 给出本书最关键的章节交接关系。

\begin{table}[htbp]
  \centering
  \footnotesize
  \caption{全书主线中的章节交接表。每一行表示“上一段向下一段交付什么”。}
  \label{tab:chapter-handoff}
  \setlength{\tabcolsep}{4pt}
  \begin{tabularx}{\textwidth}{@{}p{2.8cm}p{4.0cm}Y@{}}
    \toprule
    上游章节 & 向下游交付的核心中间量 & 下游章节如何使用 \\
    \midrule
    \cref{ch:maxwell,ch:bloch,ch:gamma} & Bloch 模、$\Gamma$ 点候选模、对称性位置、候选频段 & \cref{ch:feedback,ch:polarization,ch:bic,ch:cwt-min,ch:pwe} 用来组织辐射、偏振和降阶模型基底 \\
    \cref{ch:feedback,ch:polarization,ch:bic} & 辐射通道、偏振选择、BIC / quasi-BIC 约束 & \cref{ch:cwt-min,ch:rcwa} 用来判断哪些模值得进入阈值代理与开放散射分析 \\
    \cref{ch:cwt-min,ch:cwt-3d,ch:cmt} & 耦合矩阵、微扰灵敏度、候选排序与阈值代理语言 & \cref{ch:pwe,ch:rcwa,ch:fdtd-fem,ch:method-compare} 用来缩小全波搜索空间和解释结果变化 \\
    \cref{ch:pwe,ch:rcwa,ch:fdtd-fem,ch:method-compare} & 候选模目录、被动损耗、有限尺寸远场、方法互证结果 & \cref{ch:epitaxial-optics,ch:workflow,ch:worked-example} 用来进入器件层与完整建模链 \\
    \cref{ch:epitaxial-optics,ch:qw-gain,ch:electrical,ch:eto-coupling} & 纵向重叠、模增益映射、注入分布、温度场与工作窗口 & \cref{ch:single-mode-design,ch:robustness,ch:dynamics,ch:metrics} 用来定义真实器件层面的稳定性与性能 \\
    \cref{ch:electrical,ch:eto-coupling} & 载流子分布 $N(\rvec)$、局部温度 $T(\rvec)$ 与电热回写系数 & \cref{ch:electrical,ch:eto-coupling} 中把 $T(\rvec)$ 反馈进迁移率/增益/折射率模型，完成 L3\textrightarrow{}L4 闭环 \\
    \cref{ch:workflow,ch:worked-example,ch:failure-modes} & 正向流程、执行样例、反向诊断 & \cref{ch:device-appendix,ch:sim-appendix} 用来固化为可复核的研究模板 \\
    \bottomrule
  \end{tabularx}
\end{table}


\section{全书需要反复区分的六类边界}
到这里，可以把本章压缩成六条阅读边界：
\begin{figure}[htbp]
  \centering
  % Visual summary of six red lines
\begin{tikzpicture}[x=1cm,y=1cm,>=Latex]
  \tikzset{
    warn/.style={draw=red!70!black,fill=red!6,rounded corners,minimum width=5.0cm,minimum height=1.45cm,align=left,font=\scriptsize}
  }
  \node[warn] (r1) at (0,0) {⚠ 红线 \#1\\无限周期布洛赫模 $\neq$ 有限器件工作模};
  \node[warn,right=0.5cm of r1] (r2) {⚠ 红线 \#2\\被动冷腔共振 $\neq$ 有源阈值};
  \node[warn,below=0.45cm of r1] (r3) {⚠ 红线 \#3\\阈值增益 $g_{\mathrm{th}}$ $\neq$ 阈值电流 $I_{\mathrm{th}}$};
  \node[warn,right=0.5cm of r3] (r4) {⚠ 红线 \#4\\均匀增益/有效折射率只是代理，不是器件默认状态};
  \node[warn,below=0.45cm of r3] (r5) {⚠ 红线 \#5\\二维教学模型结论不可直接替代三维平板器件结论};
  \node[warn,right=0.5cm of r5] (r6) {⚠ 红线 \#6\\线性阈值代理 $\neq$ 阈上稳态激光理论};
\end{tikzpicture}

  \caption{“六条红线”可视化速记图。后文首次触及对应误区时，统一用“⚠ 红线 \#N”作边注提示。}
  \label{fig:ch02_six_redlines}
\end{figure}
\begin{enumerate}
  \item 无限周期的 Bloch 模不是有限器件的最终工作模（L1 \textrightarrow{} 器件层）；
  \item 被动冷腔共振不是有源阈值（L1 \textrightarrow{} L2）；
  \item 阈值增益不是阈值电流（L2 \textrightarrow{} L3/L4）；
  \item 均匀增益或有效折射率只是代理模型，不是真实器件默认状态（L2 \textrightarrow{} L3/L4）；
  \item 二维教学模型不是三维 slab 结构结论，除非已明确说明哪些纵向自由度被压缩进了等效参数（2D \textrightarrow{} 3D 层级替换）；
  \item 线性阈值代理不是阈上稳态激光理论，除非已显式引入有源自洽场或饱和非线性框架（L2 \textrightarrow{} 阈上动力学）。
\end{enumerate}
如果读者在后续章节中不断回看这六条，整本书会清楚得多。

\section*{本章小结}
本章的任务是把全书的语义基础固定下来。自此以后，看到坐标、频域符号、$\Gamma$ 点、带边、阈值增益、阈值电流等术语，都应自动联想到它们所属的模型层级与物理含义。

\section*{练习题}
\begin{enumerate}
  \item \basicq 在本书采用的时间因子约定下，解释为什么 $\tilde{\omega}=\omega_0-\ii\gamma$ 表示衰减。

  \item \midq 试着用自己的话区分“Bloch 模”“被动共振模”“激射模”。

  \item \hardq 给出一个例子，说明为什么 $g_\mathrm{th}$ 很低时，$I_\mathrm{th}$ 仍可能很高。
\end{enumerate}

\section*{建议继续阅读}
\begin{itemize}
  \item \cite{sakoda2005} 中关于晶格、倒格矢与对称性记号的章节。

  \item 建议将本书的符号表与本章一起对照阅读，随时检查各符号所属的模型层级。
\end{itemize}

